\chapter{Paths and trajectory timing}
\label{chap:lecture12}
A geometric route does not specify how fast a robot should move along it.
The same route can be followed slowly or quickly, with gradual starts and
stops or abrupt changes in the requested velocity. Motion planning must
therefore distinguish the geometry of a path from its timing. This chapter
first considers motion without obstacles, then examines the velocity and
acceleration requirements of the resulting trajectories.

\section{Angular coordinates and configuration-space plots}
Consider a planar RR arm with unrestricted revolute joints. A convenient
coordinate representation is
\[
 q_1,q_2\in(-\pi,\pi].
\]
The half-open interval avoids storing both $-\pi$ and $+\pi$ as distinct
representatives of the same angular configuration. Each angle is periodic;
crossing the right edge of a coordinate plot brings its representative
back to the left edge. The same applies to the top and bottom edges.

\begin{figure}[htbp]
\centering\begin{tikzpicture}[scale=.85]
\filldraw[fill=blue!5,thick] (-3,-3) rectangle (3,3);
\draw[axis] (-3.6,0)--(3.7,0) node[right] {$q_1$};
\draw[axis] (0,-3.6)--(0,3.7) node[above] {$q_2$};
\node[below left] at (-3,0) {$-\pi$};\node[below right] at (3,0) {$\pi$};
\node[above left] at (0,3) {$\pi$};\node[below left] at (0,-3) {$-\pi$};
\foreach \x in {-3,3} \draw[blue!70!black,very thick,->] (\x,-1.8)--(\x,-.8);
\foreach \y in {-3,3} \draw[red!70!black,very thick,->] (.8,\y)--(1.8,\y);
\fill (2.83,1) circle(.06) node[above left] {$170^\circ$};
\fill (-2.83,1) circle(.06) node[above right] {$-170^\circ$};
\draw[green!45!black,very thick,->] (2.83,1)--(3.5,1);
\draw[green!45!black,very thick,->] (-3.5,1)--(-2.83,1);
\node[align=center] at (0,-4.1) {Matching edge arrows indicate periodic identification.};
\end{tikzpicture}

\caption{A coordinate square for two unrestricted revolute joints.
Opposite edges represent matching configurations. The short motion across
the angular boundary appears as two pieces in this plot, although the
physical joint motion is continuous. The interval convention selects one
representative at each boundary.}
\label{fig:l12-periodic}
\end{figure}

The square is a representation of the periodic configuration space, not
a pair of mechanical stops. With both joints unrestricted, its topology
is a torus, $S^1\times S^1$. A joint with physical stops instead has an
interval of admissible motion, and those stops cannot be crossed merely
by wrapping the angle numerically.

For example, moving from $170^\circ$ to $-170^\circ$ can be a positive
$20^\circ$ rotation for an unrestricted joint. Directly subtracting the
stored representatives gives $-340^\circ$. Either represents a route to
the same final orientation, but they describe different motions. Choose
an unwrapped displacement consistent with the desired route and constraints,
then normalize coordinates for display if needed. Accumulated turns may
also matter when cables or other physical constraints are present.

\section{A path specifies geometry and order}
A \emph{path} describes the configurations or task values to be traversed,
including their order, without specifying the traversal times. Represent
a configuration path by
\begin{equation}
 q:[0,1]\longrightarrow\mathcal C,\qquad
 q(0)=q_S,\quad q(1)=q_G.
 \label{eq:l12-path}
\end{equation}
The parameter $\lambda\in[0,1]$ is dimensionless progress along the path;
it is not time. The set of visited points is the path's geometric image,
but a parameterized curve also records how those points are ordered.

For a segment in a chosen coordinate domain, let $\Delta q=q_G-q_S$.
Then
\begin{equation}
 \boxed{q(\lambda)=q_S+\Delta q\,\lambda
 =(1-\lambda)q_S+\lambda q_G.}
 \label{eq:l12-segment}
\end{equation}
Componentwise, $q_i(\lambda)=q_{iS}+\Delta q_i\lambda$. All joints use
the same progress parameter to trace this particular straight segment.

As an example, take
\[
 q_S=\begin{bmatrix}\pi/6\\\pi/4\end{bmatrix},\qquad
 q_G=\begin{bmatrix}\pi/2\\-\pi/2\end{bmatrix},\qquad
 \Delta q=\begin{bmatrix}\pi/3\\-3\pi/4\end{bmatrix}.
\]
The midpoint is $q(1/2)=(\pi/3,-\pi/8)\trans$. The first joint increases
while the second decreases. Both remain inside the plotted coordinate
intervals, although collision checking would still be needed if obstacles
were present.

\begin{figure}[htbp]
\centering\begin{tikzpicture}[scale=1.2]
\draw[gray] (-2,-2) rectangle (2,2);
\draw[axis] (-2.3,0)--(2.6,0) node[right] {$q_1$};
\draw[axis] (0,-2.3)--(0,2.5) node[above] {$q_2$};
\node[below] at (2,0) {$\pi$};\node[below] at (-2,0) {$-\pi$};
\node[left] at (0,2) {$\pi$};\node[left] at (0,-2) {$-\pi$};
\draw[blue!70!black,very thick,->] (.3333,.5)--(1,-1);
\fill (.3333,.5) circle(.045) node[above right] {$S=(\pi/6,\pi/4)$};
\fill (1,-1) circle(.045) node[below right] {$G=(\pi/2,-\pi/2)$};
\fill (.6667,-.25) circle(.045);
\draw (.6667,-.25)--(2.8,-.5) node[right,align=left] {$\lambda=1/2$\\$(\pi/3,-\pi/8)$};
\end{tikzpicture}

\caption{Straight interpolation in the joint coordinates for the stated
RR start and goal. The parameter $\lambda$ selects a configuration on the
segment. Its timing has not yet been assigned.}
\label{fig:l12-segment}
\end{figure}

A Cartesian position segment can be parameterized in the same way.
However, its intermediate points must be reachable, and an admissible
joint path must realize them. Orientation interpolation requires a
representation appropriate to rotations; interpolating arbitrary rotation
matrix entries does not generally preserve a valid rotation.
For a differential-drive robot, a geometric pose path must also respect
its velocity constraints. Coordinate interpolation alone does not establish
that the robot can execute a motion.

\section{A trajectory adds a timing law}
A \emph{timing law} assigns progress as a function of time:
\[
 \lambda=\lambda(t),\qquad \lambda(0)=0,\qquad\lambda(T)=1.
\]
For a forward traversal choose $0\le\lambda(t)\le1$ with
$\dot\lambda(t)\ge0$. The trajectory is the composition
\begin{equation}
 q_d(t)=q(\lambda(t)),\qquad0\le t\le T.
 \label{eq:l12-trajectory}
\end{equation}
The subscript $d$ emphasizes that this is a desired motion for a controller
to track. The actual measured configuration need not equal it exactly.

\subsection{Linear timing}
The simplest timing law is $\lambda(t)=t/T$, giving
\begin{equation}
 q_{i,d}(t)=q_{iS}+\Delta q_i\frac{t}{T}.
 \label{eq:l12-linear}
\end{equation}
For the first joint of the example, choose $T=10\,\mathrm s$:
\[
 q_{1,d}(t)=\frac\pi6+\frac\pi3\frac{t}{10\,\mathrm s}.
\]
It starts at $30^\circ$ and ends at $90^\circ$. The displacement is
$\pi/3$, or $60^\circ$; the final angle is $\pi/2$, not $\pi/3$.

\subsection{Samples for a digital controller}
At a control interval $\Delta t$, evaluate the reference at timestamps
$t_k=k\Delta t$. With $\Delta t=0.5\,\mathrm s$, the example gives
\begin{center}
\begin{tabular}{rrrr}
\toprule
$k$ & $t_k$ (s) & $\lambda(t_k)$ & $q_{1,d}(t_k)$ (degrees)\\
\midrule
0 & 0 & 0 & 30\\
1 & 0.5 & 0.05 & 33\\
2 & 1 & 0.10 & 36\\
3 & 1.5 & 0.15 & 39\\
10 & 5 & 0.50 & 60\\
20 & 10 & 1 & 90\\
\bottomrule
\end{tabular}
\end{center}
There are 21 samples including both endpoints. The $0.5\,\mathrm s$
interval is illustrative, not a recommended motor-control rate.

Software can evaluate the analytical formula directly at each control
step; symbolic algebra is not required. Alternatively, it can store a
table of reference values and timestamps. Either representation needs a
defined update and interpolation policy. If $T$ is not an integer multiple
of $\Delta t$, include the final timestamp explicitly or evaluate the
clamped final reference at subsequent updates.

The trajectory values are reference positions, not PWM commands. A
controller compares them with measured motion and computes suitable
actuator commands. Tracking error, sampling, and actuator limits remain
part of the execution problem.

\section{Why abrupt constant-speed starts are problematic}
Suppose the joint is at rest before $t=0$ and again after $t=T$. Extending
the linear reference outside the motion interval gives
\[
 q_d(t)=\begin{cases}
 q_S,&t<0,\\
 q_S+\Delta q\,t/T,&0\le t\le T,\\
 q_G,&t>T.
 \end{cases}
\]
For a scalar joint, its velocity is zero before and after the move and
$\Delta q/T$ in the interior. The jumps at the endpoints require
instantaneous changes in velocity.

\begin{figure}[htbp]
\centering\begin{tikzpicture}[x=1.25cm,y=.85cm]
\foreach \h/\lab in {6/{q_d},3/{\dot q_d},0/{\ddot q_d}} {
\draw[axis] (-.7,\h)--(6,\h) node[right] {$t$};
\draw[axis] (0,{\h-.4})--(0,{\h+2.15});
\node[left] at (-.15,{\h+2}) {$\lab$};
\node[below left] at (0,\h) {$0$};
\node[below] at (5,\h) {$T$};}
\draw[gray,dashed] (5,-1.4)--(5,8.1);
\draw[blue!70!black,very thick] (-.6,6.4)--(0,6.4)--(5,8)--(5.8,8);
\node[left] at (-.65,6.4) {$q_S$};\node[left] at (-.65,8) {$q_G$};
\draw[projection] (0,8)--(5,8);
\draw[blue!70!black,very thick] (-.6,3)--(0,3) (0,4.5)--(5,4.5) (5,3)--(5.8,3);
\draw[blue!70!black,dashed] (0,3)--(0,4.5) (5,4.5)--(5,3);
\node[left] at (0,4.5) {$\Delta q/T$};
\draw[blue!70!black,very thick] (-.6,0)--(5.8,0);
\draw[red!75!black,very thick,->] (0,0)--(0,1.5);
\draw[red!75!black,very thick,->] (5,0)--(5,-1.3);
\node[right] at (.1,1.2) {positive impulse};
\node[left] at (4.9,-1.1) {negative impulse};
\end{tikzpicture}

\caption{A linear position reference joined to stationary endpoints.
Position is continuous, but velocity jumps at $0$ and $T$. The arrows in
the acceleration plot represent ideal impulses, not finite acceleration
peaks. The sketch shows a positive displacement.}
\label{fig:l12-linear-profile}
\end{figure}

In the distributional sense, the acceleration is
\[
 \ddot q_d(t)=\frac{\Delta q}{T}\bigl[\delta(t)-\delta(t-T)\bigr].
\]
The Dirac delta is an ideal impulse with a specified integral, not an
ordinary finite-valued acceleration at one instant. A real actuator with
finite torque and a load with inertia cannot follow these velocity jumps
exactly. It needs time to accelerate and brake. The resulting tracking
behavior depends on the controller and system dynamics; overshoot is a
possibility, not an unavoidable outcome for every implementation.

Available acceleration depends on torque, inertia, loading, transmission,
and speed-dependent limits. Motor power alone does not specify a constant
joint-acceleration bound. For trajectory design, a simplified model can
nevertheless use conservative velocity and acceleration limits.

\section{Acceleration, cruising, and braking}
\label{sec:l12-trapezoid}
A common rest-to-rest motion uses constant acceleration, an optional
constant-speed interval, and constant deceleration. Integrating these
three phases produces a trapezoidal velocity profile and a position
profile that is quadratic, then linear, then quadratic.

\begin{figure}[htbp]
\centering\begin{tikzpicture}[x=1.3cm,y=.85cm]
\foreach \h/\lab in {6/{q_d},3/{\dot q_d},0/{\ddot q_d}} {
\draw[axis] (-.6,\h)--(5.8,\h) node[right] {$t$};
\draw[axis] (0,{\h-.3})--(0,{\h+2.15});
\node[left] at (-.15,{\h+2}) {$\lab$};
\foreach \x/\s in {1/{t_1},4/{t_2},5/{T}} \node[below] at (\x,\h) {$\s$};}
\foreach \x in {1,4,5} \draw[gray,dashed] (\x,-1.3)--(\x,8);
\draw[blue!70!black,very thick] (-.5,6.3)--(0,6.3);
\draw[blue!70!black,very thick,domain=0:1,samples=20] plot (\x,{6.3+.2125*\x*\x});
\draw[blue!70!black,very thick] (1,6.5125)--(4,7.7875);
\draw[blue!70!black,very thick,domain=4:5,samples=20] plot (\x,{8-.2125*(5-\x)*(5-\x)});
\draw[blue!70!black,very thick] (5,8)--(5.6,8);
\node[left] at (-.65,6.3) {$q_S$};\node[left] at (-.65,8) {$q_G$};
\draw[blue!70!black,very thick] (-.5,3)--(0,3)--(1,4.5)--(4,4.5)--(5,3)--(5.6,3);
\node[left] at (0,4.5) {$v_p$};\draw[projection] (0,4.5)--(1,4.5);
\draw[blue!70!black,very thick] (-.5,0)--(0,0) (0,1.1)--(1,1.1) (1,0)--(4,0) (4,-1.1)--(5,-1.1) (5,0)--(5.6,0);
\draw[blue!70!black,dashed] (0,0)--(0,1.1) (1,1.1)--(1,0) (4,0)--(4,-1.1) (5,-1.1)--(5,0);
\node[left] at (0,1.1) {$+A$};\node[left] at (0,-1.1) {$-B$};
\end{tikzpicture}

\caption{A rest-to-rest profile with acceleration, cruising, and braking.
Velocity is continuous; acceleration switches between finite levels.
Position and velocity join continuously at $t_1$ and $t_2$. The acceleration
and deceleration magnitudes need not be equal in general.}
\label{fig:l12-trapezoid}
\end{figure}

\subsection{A scalar construction}
The following formulas make the three-phase construction explicit. Let
$D=|q_G-q_S|>0$ and $\sigma=\operatorname{sign}(q_G-q_S)$.
Write $q_d(t)=q_S+\sigma u(t)$, where progress $u$ increases from 0 to $D$.
Let $A>0$ be the acceleration limit along the move, $B>0$ the braking
magnitude, and $V>0$ the speed limit. Suppose the peak speed is $v_p$.
The acceleration and braking times are
\[
 t_a=\frac{v_p}{A},\qquad t_d=\frac{v_p}{B},
\]
and the distances used by those phases are $v_p^2/(2A)$ and $v_p^2/(2B)$.
If the speed limit can be reached, set $v_p=V$ and
\begin{equation}
 t_c=\frac{D-V^2/(2A)-V^2/(2B)}{V}\ge0,
 \quad t_1=t_a,\quad t_2=t_a+t_c,\quad T=t_a+t_c+t_d.
 \label{eq:l12-times}
\end{equation}
The resulting progress is
\begin{equation}
 u(t)=\begin{cases}
 \tfrac12At^2,&0\le t\le t_1,\\
 \dfrac{v_p^2}{2A}+v_p(t-t_1),&t_1\le t\le t_2,\\
 D-\tfrac12B(T-t)^2,&t_2\le t\le T.
 \end{cases}
 \label{eq:l12-piecewise}
\end{equation}
The overlapping endpoint definitions agree. The associated timing law is
$\lambda(t)=u(t)/D$. The velocity ramps from zero to $v_p$, remains at
$v_p$ during cruising, and returns to zero during braking.

\subsection{Short moves: triangular velocity}
If $D<V^2/(2A)+V^2/(2B)$, there is insufficient distance to reach $V$
and still stop at the target. Set $t_c=0$ and solve for the lower peak speed:
\begin{equation}
 v_p=\sqrt{\frac{2DAB}{A+B}}<V.
 \label{eq:l12-triangle}
\end{equation}
The velocity profile is triangular. Acceleration switches directly from
$+A$ to $-B$; this is the two-phase bang-bang case. The trapezoidal case
adds a zero-acceleration cruise interval and is often called
bang-coast-bang. Both are idealized profiles with jumps in acceleration;
a jerk-limited design would smooth those changes further.

These profiles give minimum-time rest-to-rest scalar motion under the
stated constant speed and acceleration bounds, without a jerk bound or
additional constraints. This does not establish global time optimality
for a coupled robot with configuration-dependent torque limits.

\subsection{Worked example}
Move from $30^\circ$ to $90^\circ$, so $D=60^\circ$. For easy numerical
interpretation use degrees here, with $V=15^\circ/\mathrm s$ and
$A=B=15^\circ/\mathrm s^2$. Equivalently, each limit has numerical
factor $\pi/12$ in radian units. Then
\[
 t_a=t_d=1\,\mathrm s,\qquad
 t_c=\frac{60-7.5-7.5}{15}=3\,\mathrm s,\qquad T=5\,\mathrm s.
\]
With $t$ measured in seconds and $q_d$ in degrees,
\[
 q_d(t)=\begin{cases}
 30+7.5t^2,&0\le t\le1,\\
 37.5+15(t-1),&1\le t\le4,\\
 90-7.5(5-t)^2,&4\le t\le5.
 \end{cases}
\]
The joint reaches $37.5^\circ$ at 1 s, begins braking at $82.5^\circ$
at 4 s, and reaches $90^\circ$ at rest at 5 s. The area under the velocity
curve is $7.5+45+7.5=60^\circ$, as required. This is a different timing
choice from the earlier 10-second linear example, for the same endpoints.

\section{Coordinating several joints}
For the straight joint path, a shared timing law gives
\begin{equation}
 \dot q_{i,d}=\Delta q_i\dot\lambda,\qquad
 \ddot q_{i,d}=\Delta q_i\ddot\lambda.
 \label{eq:l12-shared}
\end{equation}
All joints finish together and preserve the selected configuration-space
segment. Under symmetric individual speed bounds $V_i$ and acceleration
magnitude bounds $A_i$, a sufficient choice is
\[
 |\dot\lambda|\le\min_{\Delta q_i\ne0}\frac{V_i}{|\Delta q_i|},\qquad
 |\ddot\lambda|\le\min_{\Delta q_i\ne0}\frac{A_i}{|\Delta q_i|}.
\]
A joint with $\Delta q_i=0$ remains fixed and is omitted from these ratios.
Mixed revolute and prismatic units are consistent because each bound is
divided by that joint's own displacement. If all displacements are zero,
the reference simply remains at the starting configuration.

Assigning unrelated timing laws to separate joints generally changes the
path through configuration space, even if every joint eventually reaches
its target. For a curved path the derivatives also involve the path's
geometry, so the simple expressions above no longer suffice.

Geometric collision checking and timing design address different parts
of the motion problem. A smooth timing law cannot repair a path that
collides with a static obstacle. Conversely, a collision-free geometric
path still needs feasible timing and a controller that can track it.

\clearpage
\section{Additional exercises}
\begin{enumerate}
\item \textbf{Angle wrapping.} An unrestricted revolute joint starts at
$170^\circ$ and must finish at $-170^\circ$. Write both the direct
coordinate difference and the short positive displacement. Explain why
the latter is not permitted if mechanical stops are at $\pm175^\circ$.

\item \textbf{A joint-space segment.} For the RR start and goal in this
chapter, compute $q(0)$, $q(1/2)$, and $q(1)$. Explain why a straight
joint-space segment need not produce a straight tool path.

\item \textbf{Sampled reference.} For $q_{1,d}(t)=\pi/6+(\pi/3)t/10$
with $t$ in seconds, calculate the references at $t=2$, 5, and 8 s. How
many samples include both endpoints when $\Delta t=0.5\,\mathrm s$?
Distinguish these reference values from actuator commands.

\item \textbf{Derivative interpretation.} Sketch position, velocity,
and acceleration for a negative displacement with linear timing and
stationary endpoints. State the signs and strengths of the ideal
acceleration impulses.

\item \textbf{Trapezoidal motion.} Reproduce the $30^\circ$ to $90^\circ$
worked example. Verify position and velocity continuity at 1 s and 4 s,
and confirm the traveled angle from the velocity-curve area.

\item \textbf{A short move.} Keep $V=15^\circ/\mathrm s$ and
$A=B=15^\circ/\mathrm s^2$, but reduce the displacement to $10^\circ$.
Determine whether cruising occurs, then compute the peak speed and total
duration. Check that the velocity-curve area equals the displacement.

\item \textbf{Shared timing.} Two joints move by $\pi/3$ radians and
$0.4\,\mathrm m$. Their speed bounds are $\pi/6\,\mathrm{rad/s}$ and
$0.1\,\mathrm{m/s}$. Find the bound on $\dot\lambda$ for a shared
straight path. Which joint sets it? Explain why this speed calculation
alone does not ensure rest-to-rest feasibility.
\end{enumerate}
